\chapter{Charts}
\label{ch:charts}

The \tab{Chart} tab plots values over time while a test runs: map registers
by name, or addresses of a manual read.

\screenshot[8cm]{chart}{Voltages and currents of a meter, one chart per unit.}
  {The Chart tab with a map: two charts (V and A) with three lines each, the
   legend with the newest values, and the hover tip showing values at one
   moment.}

\section{Adding lines}

\begin{itemize}
  \item With a map: \field{Add a register to the chart} lists the map's
        numeric registers.
  \item With a manual read: add an address from the list, or right-click a
        cell on the \tab{Read registers} tab and choose \nav{Add to chart}.
\end{itemize}

Up to eight lines can be shown, each in its own colour. Lines with the same
unit share a chart, so values of different sizes (volts and amperes) do not
flatten each other. The tab's name counts the lines: \tab{Chart (3)}.

Values are recorded from the moment a line is added, while the test runs. To
remove a line, click \ui{×} next to its name in the legend.

\section{Reading the chart}

\begin{description}[style=nextline, leftmargin=0pt, font=\intermedium]
  \item[Time span] \field{Time span} shows the last minute, 5 or 15 minutes,
    the last hour, or everything since the lines were added.
  \item[Legend] Each line's name and its newest value.
  \item[Hover] Move the pointer over a chart to see every line's value at that
    moment.
  \item[Spikes] When there are more samples than pixels, each pixel keeps its
    lowest and highest value, so short spikes stay visible.
  \item[Gaps] A failed read leaves a gap in the line instead of a false value.
\end{description}

\ui{Clear} empties the charts and keeps the lines.

\section{Saving a chart as an image}
\label{sec:chart-image}

\ui{Export image…} copies or saves the charts shown, for a report or a bug
report:

\begin{center}
\begin{tabularx}{\linewidth}{l L}
  \toprule
  \thead{Choice} & \thead{What you get} \\
  \midrule
  Copy image & The same picture as PNG, on the clipboard, to paste into a
    ticket, an e-mail or a document. \ui{Copied.} says it worked. \\
  Copy image, light theme & As Copy image, in the light theme. \\
  PNG… & A picture at twice the screen size, sharp in documents. \\
  SVG… & A vector drawing that scales to any size and opens in other
    programs (colours written out, no styles needed). \\
  PNG, light theme… & As PNG, in the light theme even when Modbench is
    dark. \\
  SVG, light theme… & As SVG, in the light theme. \\
  \bottomrule
\end{tabularx}
\end{center}

The image holds every chart shown, each with its unit and legend, under a line
with the date and the time span. The hover marks are left out. The image is
taken when you choose, so it does not move on while the save dialog is open.

The Trend player (Section~\ref{sec:trend-player}) has the same button.
